% Options for packages loaded elsewhere
\PassOptionsToPackage{unicode}{hyperref}
\PassOptionsToPackage{hyphens}{url}
\PassOptionsToPackage{dvipsnames,svgnames,x11names}{xcolor}
\documentclass[
]{article}
\usepackage{xcolor}
\usepackage[margin=1in]{geometry}
\usepackage{amsmath,amssymb}
\setcounter{secnumdepth}{-\maxdimen} % remove section numbering
\usepackage{iftex}
\ifPDFTeX
  \usepackage[T1]{fontenc}
  \usepackage[utf8]{inputenc}
  \usepackage{textcomp} % provide euro and other symbols
\else % if luatex or xetex
  \usepackage{unicode-math} % this also loads fontspec
  \defaultfontfeatures{Scale=MatchLowercase}
  \defaultfontfeatures[\rmfamily]{Ligatures=TeX,Scale=1}
\fi
\usepackage{lmodern}
\ifPDFTeX\else
  % xetex/luatex font selection
  \setmainfont[]{Helvetica Neue}
\fi
% Use upquote if available, for straight quotes in verbatim environments
\IfFileExists{upquote.sty}{\usepackage{upquote}}{}
\IfFileExists{microtype.sty}{% use microtype if available
  \usepackage[]{microtype}
  \UseMicrotypeSet[protrusion]{basicmath} % disable protrusion for tt fonts
}{}
\makeatletter
\@ifundefined{KOMAClassName}{% if non-KOMA class
  \IfFileExists{parskip.sty}{%
    \usepackage{parskip}
  }{% else
    \setlength{\parindent}{0pt}
    \setlength{\parskip}{6pt plus 2pt minus 1pt}}
}{% if KOMA class
  \KOMAoptions{parskip=half}}
\makeatother
\usepackage{longtable,booktabs,array}
\usepackage{caption}
\captionsetup[table]{skip=6pt}
\newcounter{none} % for unnumbered tables
\usepackage{calc} % for calculating minipage widths
% Correct order of tables after \paragraph or \subparagraph
\usepackage{etoolbox}
\makeatletter
\patchcmd\longtable{\par}{\if@noskipsec\mbox{}\fi\par}{}{}
\makeatother
\setlength{\emergencystretch}{3em} % prevent overfull lines
\providecommand{\tightlist}{%
  \setlength{\itemsep}{0pt}\setlength{\parskip}{0pt}}
\usepackage{bookmark}
\IfFileExists{xurl.sty}{\usepackage{xurl}}{} % add URL line breaks if available
\urlstyle{same}
% fallback for those not using the hyperref driver hyperxmp:
\makeatletter
\@ifundefined{xmpquote}{\newcommand{\xmpquote}[1]{#1}}{}
\makeatother
\hypersetup{
  pdftitle={Reporte de Avance No. 6},
  pdfauthor={Equipo de Desarrollo ABAQ},
  colorlinks=true,
  linkcolor={Maroon},
  filecolor={Maroon},
  citecolor={Blue},
  urlcolor={Blue},
  pdfcreator={LaTeX via pandoc}}

\title{Reporte de Avance No. 6}
\author{Equipo de Desarrollo ABAQ}
\date{}

\begin{document}
\maketitle

{
\setcounter{tocdepth}{3}
\tableofcontents
}
\section{Reporte de Avance No.~6}\label{reporte-de-avance-no.-6}

{\def\LTcaptype{none} % do not increment counter
\begin{longtable}[]{@{}ll@{}}
\toprule\noalign{}
\endhead
\bottomrule\noalign{}
\endlastfoot
\textbf{Fecha} & 26 de octubre de 2026 \\
\textbf{Versión} & v1.1.0 \\
\textbf{Estado} & En desarrollo (Infraestructura y Despliegue) \\
\textbf{Preparado por} & Equipo de Desarrollo ABAQ y Consultoría
Técnica \\
\end{longtable}
}

\begin{center}\rule{0.5\linewidth}{0.5pt}\end{center}

\textbf{Descargas} Puedes descargar este documento en los siguientes
formatos:

\begin{verbatim}
-  **[Descargar PDF](latex/report6.pdf)**
-  **[Descargar Word (.docx)](latex/report6.docx)**
-  **[Descargar fuente LaTeX (.tex)](latex/report6.tex)**
-  **[Ver grabación de la sesión](https://teams.microsoft.com/l/meetingrecap?driveId=b%2107Fn-TnF5Ui0l8Rgu8qBvlEeGZaatelEubk4irn-v_aq61IqHa-IQoy4hgnYmHud&driveItemId=01VBXFYXFSLVC5KGW6LFCKT7JYHLYCA7JK&sitePath=https%3A%2F%2Fabaq2-my.sharepoint.com%2F%3Av%3A%2Fg%2Fpersonal%2Fmar_ledezma_abaq_org_mx%2FIQCyXUXVGt5ZRKn9ODrwIH0qAQiXWSgAp1MfhJG8hoJFHEY&fileUrl=https%3A%2F%2Fabaq2-my.sharepoint.com%2Fpersonal%2Fmar_ledezma_abaq_org_mx%2FDocuments%2FGrabaciones%2FAvances+IT+ABAQ-20261026_91334%E2%80%AFa.m.-Grabaci%C3%B3n+de+la+reuni%C3%B3n.mp4%3Fweb%3D1&iCalUid=040000008200E00074C5B7101A82E00800000000F8C496A6EB64DD010000000000000000100000001303FA33EEEBC24483F036F8FF6E5D1E&threadId=19%3Ameeting_N2ZhODk3ODUtNmQ1Ny00YjRlLWJlMTYtMjc4YmQzMzYzZGEx%40thread.v2&organizerId=eec797ff-0209-48ac-92ea-b2c078b62786&tenantId=2876b5c1-692f-4d5a-ba3c-e4958f8791ce&callId=33a4accd-2688-4fc8-b14d-1b4153ad36b6&threadType=meeting&meetingType=Scheduled&subType=RecapSharingLink_RecapCore&recapType=RecordingAndTranscript)**
\end{verbatim}

\subsection{Resumen ejecutivo}\label{resumen-ejecutivo}

Este documento presenta el sexto reporte de avance del proyecto ABAQ.
Durante esta sesión de trabajo se definió el plan de acción, la
arquitectura de infraestructura en la nube y los proveedores de
servicios que se utilizarán para el despliegue y puesta en producción de
los dos aplicativos web de la organización: el \textbf{Sistema de Base
de Datos} y el \textbf{Sistema de Gestión de Tareas}.

En la reunión participaron Ángel, Fernando Israel Ríos García
(voluntario del equipo de base de datos y documentación), y los
consultores profesionales Michelle y Leo, quienes asumirán la
administración de la infraestructura técnica y los entornos de nube.
Durante la sesión se otorgaron los accesos requeridos a la base de datos
MongoDB y a la organización de GitHub donde residen los repositorios
oficiales de ambos proyectos.

Asimismo, se acordó la política de direccionamiento web mediante la
adopción del dominio de segundo nivel con TLD \texttt{.site}
(\texttt{nombreAplicacion.abaq.site}), preservando de forma íntegra el
dominio institucional principal (\texttt{abaq.org.mx}). Para el cómputo
en la nube se seleccionó \textbf{Railway} (\texttt{railway.app}) como
plataforma homogénea de despliegue. En este período no se registraron
modificaciones de código en backend ni frontend, centrándose la
totalidad del esfuerzo en los acuerdos estratégicos de despliegue y
soporte operativo.

\begin{center}\rule{0.5\linewidth}{0.5pt}\end{center}

\subsection{1. Introducción}\label{introducciuxf3n}

El presente reporte documenta los acuerdos establecidos en la sexta
sesión oficial de seguimiento (26 de octubre de 2026), cuyo objetivo
primordial fue coordinar la infraestructura de despliegue para las
aplicaciones desarrolladas por el equipo de servicio social durante el
período agosto--diciembre de 2026.

Los temas centrales abordados en la sesión comprenden:

\begin{itemize}
\tightlist
\item
  \textbf{Equipo de trabajo e integración de consultores:} Coordinación
  entre los desarrolladores del sistema y el equipo de consultoría
  técnica (Michelle y Leo) para la futura administración de servidores.
\item
  \textbf{Acceso a repositorios y base de datos:} Transferencia de
  accesos al código fuente en GitHub y a la base de datos central en
  MongoDB.
\item
  \textbf{Estrategia de dominios y resolución DNS:} Selección del
  dominio \texttt{abaq.site} y estructura de subdominios para aislar los
  entornos internos de la imagen institucional.
\item
  \textbf{Infraestructura en la nube:} Selección de Railway como
  proveedor de cómputo en la nube para el alojamiento unificado de
  servicios y bases de datos.
\end{itemize}

\begin{center}\rule{0.5\linewidth}{0.5pt}\end{center}

\subsection{2. Acuerdos y Decisiones de
Infraestructura}\label{acuerdos-y-decisiones-de-infraestructura}

\subsubsection{2.1 Integración del equipo técnico y
consultoría}\label{integraciuxf3n-del-equipo-tuxe9cnico-y-consultoruxeda}

Se formalizó la incorporación de Michelle y Leo como consultores
profesionales y administradores de la infraestructura técnica de ABAQ.
Su función principal consiste en gestionar el aprovisionamiento de
servidores en la nube y coordinar la configuración de los dominios
institucionales.

{\def\LTcaptype{none} % do not increment counter
\begin{longtable}[]{@{}
  >{\raggedright\arraybackslash}p{(\linewidth - 2\tabcolsep) * \real{0.5000}}
  >{\raggedright\arraybackslash}p{(\linewidth - 2\tabcolsep) * \real{0.5000}}@{}}
\toprule\noalign{}
\begin{minipage}[b]{\linewidth}\raggedright
Participante
\end{minipage} & \begin{minipage}[b]{\linewidth}\raggedright
Rol / Función
\end{minipage} \\
\midrule\noalign{}
\endhead
\bottomrule\noalign{}
\endlastfoot
Ángel & Coordinación y seguimiento de proyectos \\
Fernando Israel Ríos García & Voluntario ABAQ --- Base de datos,
documentación y gestión técnica de DNS \\
Johan Del Ángel González & Desarrollador principal --- Arquitectura y
soporte de base de datos \\
Michelle & Consultora profesional --- Administración de infraestructura
y nube \\
Leo & Consultor profesional --- Administración de infraestructura y
nube \\
\end{longtable}
}

\subsubsection{2.2 Accesos y repositorios del
proyecto}\label{accesos-y-repositorios-del-proyecto}

Durante la sesión se proporcionaron los siguientes accesos al equipo de
consultoría:

\begin{enumerate}
\def\labelenumi{\arabic{enumi}.}
\tightlist
\item
  \textbf{Base de Datos MongoDB:} Acceso administrativo al clúster de
  base de datos del aplicativo de voluntarios.
\item
  \textbf{Organización en GitHub:} Incorporación a la organización
  oficial donde se encuentran resguardados los proyectos de software:

  \begin{itemize}
  \tightlist
  \item
    Repositorio del proyecto \textbf{Base de Datos ABAQ}.
  \item
    Repositorio del proyecto \textbf{Gestión de Tareas} (incorporado el
    27 de octubre de 2026).
  \end{itemize}
\item
  \textbf{Inducción técnica:} Se realizó una sesión de contexto técnico
  explicando el alcance, modelo de datos y arquitectura de ambos
  sistemas para resolver dudas iniciales de los administradores.
\end{enumerate}

\subsubsection{2.3 Estrategia de dominio y administración
DNS}\label{estrategia-de-dominio-y-administraciuxf3n-dns}

Para garantizar una operación segura y profesional, se determinaron los
siguientes lineamientos sobre las direcciones web:

\begin{itemize}
\tightlist
\item
  \textbf{Preservación del dominio institucional:} El dominio principal
  (\texttt{abaq.org.mx}) no se utilizará para los aplicativos internos
  de gestión, resguardando la marca e identidad corporativa de la
  asociación.
\item
  \textbf{Uso del dominio \texttt{.site}:} Se determinó emplear el TLD
  \texttt{.site}, adoptando subdominios específicos bajo el formato:
  \texttt{nombreAplicacion.abaq.site} (por ejemplo,
  \texttt{basededatos.abaq.site} y \texttt{tareas.abaq.site}).
\item
  \textbf{Panel y gestión de registros DNS:} Debido a las restricciones
  técnicas del proveedor actual (\texttt{Akky.mx}), que no cuenta con un
  gestor ágil de zonas DNS, se acordó definir un nuevo proveedor de
  gestión de dominio. En este nuevo panel se otorgará un rol técnico a
  Fernando Israel Ríos García con facultades para crear y actualizar los
  registros A, CNAME e IPs de los servicios.
\end{itemize}

\subsubsection{2.4 Plataforma de despliegue en la nube
(Railway)}\label{plataforma-de-despliegue-en-la-nube-railway}

Se seleccionó \textbf{Railway (\texttt{railway.app})} como el proveedor
oficial de Infraestructura como Servicio (IaaS/PaaS) para ABAQ por las
siguientes ventajas:

\begin{itemize}
\tightlist
\item
  Despliegue homogéneo y automatizado desde los repositorios de GitHub.
\item
  Soporte nativo para contenedores y servicios desacoplados (Frontend
  Angular / React, Backend Node.js y bases de datos).
\item
  Gestión simplificada de variables de entorno, certificados SSL
  automáticos y métricas de consumo.
\end{itemize}

\begin{center}\rule{0.5\linewidth}{0.5pt}\end{center}

\subsection{3. Cambios en el Servidor
(Backend)}\label{cambios-en-el-servidor-backend}

Sin modificaciones registradas en código fuente durante este período. La
sesión estuvo orientada exclusivamente a la definición de
infraestructura y acuerdos con los consultores técnicos.

\begin{center}\rule{0.5\linewidth}{0.5pt}\end{center}

\subsection{4. Cambios en la Aplicación Web
(Frontend)}\label{cambios-en-la-aplicaciuxf3n-web-frontend}

Sin modificaciones registradas en código fuente durante este período.

\begin{center}\rule{0.5\linewidth}{0.5pt}\end{center}

\subsection{5. Evidencia Visual}\label{evidencia-visual}

No aplica para esta entrega, debido a que no hubo modificaciones de
interfaz ni nuevas pantallas desplegadas. La evidencia de los acuerdos
tomados se encuentra registrada en la grabación oficial de la sesión en
Microsoft Teams vinculada al inicio de este documento.

\begin{center}\rule{0.5\linewidth}{0.5pt}\end{center}

\subsection{6. Estado Actual del
Proyecto}\label{estado-actual-del-proyecto}

\textbf{Fase de Infraestructura y Despliegue} Se consolidó la estrategia
de infraestructura y se incorporó al equipo de consultoría técnica
externa. Los aplicativos de Base de Datos y Gestión de Tareas se
encuentran en preparación para su publicación en la nube bajo el dominio
\texttt{abaq.site}.

\begin{center}\rule{0.5\linewidth}{0.5pt}\end{center}

\subsection{7. Próximos Pasos}\label{pruxf3ximos-pasos}

\begin{enumerate}
\def\labelenumi{\arabic{enumi}.}
\tightlist
\item
  \textbf{Infraestructura en Railway:}

  \begin{itemize}
  \tightlist
  \item
    1.1 Crear la cuenta institucional de Railway para ABAQ, documentar
    las credenciales y compartirlas de manera segura con el equipo de
    consultores.
  \item
    1.2 Configurar el entorno de proyectos y enlazar los repositorios de
    GitHub correspondientes.
  \end{itemize}
\item
  \textbf{Gestión de Dominio y DNS:}

  \begin{itemize}
  \tightlist
  \item
    2.1 Definir los identificadores definitivos de subdominio para la
    aplicación de Base de Datos y la aplicación de Gestión de Tareas.
  \item
    2.2 Concretar la adquisición o migración del dominio
    \texttt{abaq.site} al nuevo proveedor de DNS.
  \item
    2.3 Asignar el perfil técnico a Fernando Israel Ríos García para la
    configuración de registros DNS hacia las IPs de Railway.
  \end{itemize}
\item
  \textbf{Canal de Soporte y Acompañamiento Técnico:}

  \begin{itemize}
  \tightlist
  \item
    3.1 Establecer el canal de atención para dudas y consultas
    operativas vía llamada telefónica o WhatsApp, en un horario de lunes
    a viernes de 09:00 a 18:00 horas (CST).
  \item
    3.2 Para dudas inherentes a la arquitectura del sistema de base de
    datos, canalizar la comunicación con Johan Del Ángel González
    (desarrollador principal) y Fernando Israel Ríos García.
  \end{itemize}
\end{enumerate}

\end{document}
